\begin{center}{\Large\bfseries Supplementary material}\\[0.5em] When Husbands Stop Working, Do Wives Start? Evidence from India's Urban Labour Force Panel\end{center}

\section*{S1. Panel linkage and sample construction}
The urban PLFS visits each household in four consecutive quarters. The 2022--23 files label July--September 2022 as quarter 5 (the panel's running quarter), and the 2023--24 files label July--September 2023 as quarter 1. I map both to a calendar index running from 1 (July--September 2022) to 8 (April--June 2024). For a record at visit $k$ observed in calendar quarter $c$, the first visit took place in quarter $c-k+1$. A person is identified by state, first-stage unit, sub-block, second-stage stratum, household and person serial numbers, together with the quarter of the first visit. The person identifiers are reused across rotation groups, and adding the first-visit quarter makes each identifier unique: no person appears twice at the same visit. In 2017--18 the revisit files label quarters one step ahead of the first-visit file; empirical matching shows that revisit $k$ falls $k$ labelled quarters after the first visit, and I shift the labels accordingly. Revisits whose first visit falls before the period covered by the files have no first-visit weights or household information and are dropped. Table~S1 reports the counts.

\begin{table}[htbp]\centering\caption*{Table S1: Panel linkage and sample construction}\small
%%TABLES1%%
\end{table}

\section*{S2. Validation against published estimates}
Table~S2 compares weighted first-visit CWS labour force participation rates (LFPR) and worker-population ratios (WPR) for persons aged 15 and above with the MoSPI press note of 23 September 2024 (Tables 4 and 5). All 24 values match to one decimal place.

\begin{table}[htbp]\centering\caption*{Table S2: Reproduced and published CWS estimates, persons aged 15 and above}\small
%%TABLES2%%
\end{table}

\section*{S3. Robustness, event-study coefficients, attrition and husbands' destinations}
\begin{table}[htbp]\centering\caption*{Table S3: Robustness and supporting estimates (percentage points)}\scriptsize
%%TABLES3%%
\par\vspace{4pt}\begin{minipage}{0.95\linewidth}\footnotesize SE: standard error. Panel A: sample and controls as in Table 2, column 1, unless stated. The two subperiod rows sum to 2 fewer transitions than the baseline because one singleton fixed-effect observation is dropped in each subperiod regression. Panel B: Sun--Abraham estimates plotted in Figure 1. Panel C: couples observed at visits 1--3 whose next visit falls within the data window; the difference is from a regression on husband's status with state-by-quarter and visit fixed effects. The second difference restricts to couples with the wife out of the labour force (the entry sample's population). The attrition sensitivity bounds add these extra leavers back to the stop group, assuming all of them (upper) or none of them (lower) entered the labour force, and shift the entry estimate by the change in the stop group's entry rate; attrition among the other couples is ignored, and the last row repeats the upper bound at the two ends of the gap's 95\% confidence interval; with stop-group entry rate among stayers $p_1=3.93$ per cent, attrition $a_1=2.19$ per cent among couples with the husband employed and the wife out of the labour force, retention $r_1=1-a_1$ and gap $g=1.13$ points, the stop group's entry rate becomes $[p_1(r_1-g)+g]/r_1$ if all extra leavers entered and $p_1(r_1-g)/r_1$ if none did, a shift of $+1.11$ and $-0.05$ points; adding the $+1.11$ shift to the baseline confidence interval's upper limit of 3.97 gives about 5.1, and the shift at the gap's upper limit (1.89) gives about 5.9; the bounds are asymmetric because the entry rate among stayers after a stop is low (3.93 per cent). Panel D: weighted shares of husband stops in the main sample.\end{minipage}
\end{table}

\section*{S4. Why the calendar-year 2024 file is excluded}
The PLFS calendar-year 2024 release (its person and household first-visit files, CPERV1 and CHHV1) contains only first-visit records for panel 4, quarters 3 to 6. Its quarters 3 and 4 are the same persons, with the same identifiers, as the 2023--24 first visits in those quarters (Table~S4). Quarters 5 and 6 are first visits whose revisits are not in the release. The file therefore adds no quarter-to-quarter transitions, and using it would double-count first visits already in the 2023--24 data. It is not used.

\begin{table}[htbp]\centering\caption*{Table S4: Urban person records in the calendar-year 2024 first-visit file}\small
%%TABLES4%%
\par\vspace{4pt}\begin{minipage}{0.95\linewidth}\footnotesize V1: first visit. Quarters are numbered as in the PLFS panel files.\end{minipage}
\end{table}
